\section{Giới Thiệu}
\label{sec:intro_vi}

Việc triển khai container dùng chung một nhân Linux khiến giao diện
lời gọi hệ thống (syscall) trở thành bề mặt tấn công có giá trị cao.
Framework DeSFAM~\cite{desfam2025} đề xuất \emph{SyscallAD}, bộ phát
hiện bất thường lai gồm Bộ Mã Hóa Biến Phân (VAE) và Isolation Forest
(iForest) trên đặc trưng thủ công của các cửa sổ syscall cố định độ
dài, triển khai dòng (inline) qua eBPF/Tetragon~\cite{tetragon}. Trên
bộ dữ liệu DongTing~\cite{dongting}, nhóm tác giả gốc công bố
$F_{1}=0{,}92$, độ chính xác $0{,}94$ và độ thu hồi $0{,}90$ với độ
trễ phát hiện dưới mili-giây.

Hai đặc tính của kết quả đó là động lực cho một phép tái hiện độc
lập. Thứ nhất, mô tả công bố của bước hợp nhất điểm số bỏ qua bộ
chuẩn hoá áp dụng lên hai kênh điểm thành phần --- một chi tiết hoá
ra quyết định việc tập hợp được hợp nhất hay sụp đổ. Thứ hai, lớp lành
tính của DongTing được lấy từ các bộ kiểm thử nhân Linux (LTP,
Kselftests, OpenPOSIX, glibc) và lớp tấn công từ các PoC của
syzkaller~\cite{syzkaller}; cả hai phân phối đều không đại diện cho
lưu lượng syscall do các khối công việc container thực sinh ra. Do đó
việc các chỉ số công bố có tổng quát hoá sang triển khai cloud-native
thực tế hay không là một câu hỏi mở.

Báo cáo này trả lời các câu hỏi nghiên cứu sau:
\begin{itemize}
  \item[\textbf{CH1.}] Bộ phát hiện SyscallAD có tái hiện được trên
        DongTing với cấu hình đã công bố hay không?
  \item[\textbf{CH2.}] Một mô hình SyscallAD huấn luyện trên DongTing
        có chuyển giao được tới vết syscall thu từ các khối công việc
        container kiểu sản xuất hay không?
  \item[\textbf{CH3.}] Các bộ phân loại chuỗi có-giám-sát (LSTM,
        1-D~CNN) so sánh như thế nào với tập hợp VAE + iForest
        không-giám-sát trên cùng không gian đặc trưng?
  \item[\textbf{CH4.}] Pipeline triển khai đạt độ trễ phát hiện nào
        trước các vết khai thác CVE thật và các hành vi tấn công phổ
        biến, đo end-to-end qua Tetragon?
\end{itemize}

Các đóng góp của báo cáo: (i)~một bản tái hiện mã nguồn mở của
SyscallAD kèm bản sửa lỗi chuẩn hoá đã được ghi chép trong pipeline
tham chiếu (Mục~\ref{sec:reproduction_vi}); (ii)~một bộ phát hiện
nguyên gRPC tiêu thụ sự kiện Cilium Tetragon ở $9$--$13$~ms cho mỗi
cửa sổ $200$ syscall trên một lõi CPU đơn (Mục~\ref{sec:deployment_vi});
(iii)~một tập dữ liệu syscall có nhãn thu trực tiếp từ các khối công
việc container kiểu sản xuất trong Kubernetes, cùng với hai mô hình
huấn luyện lại (Mục~\ref{sec:methodology_vi}); và (iv)~một phép đánh
giá so sánh năm biến thể mô hình trên ba tập kiểm tra kèm độ trễ phát
hiện theo từng CVE (Mục~\ref{sec:evaluation_vi}).

\FloatBarrier
\section{Bối Cảnh và Công Trình Liên Quan}
\label{sec:background_vi}

\subsection{Phát hiện bất thường dựa trên syscall}
Forrest và cộng sự~\cite{forrest1996} đã giới thiệu chuỗi syscall như
một dấu vân hành vi quá trình. Các công trình tiếp theo khai thác mô
hình chuỗi (LSTM/CNN) và phương pháp không-giám-sát như autoencoder và
isolation forest~\cite{ref38_nimos,ref27_optimus}. Nhóm tác giả DeSFAM
kết hợp hai phương pháp sau thành tập hợp cộng, với trực giác rằng
tái tạo VAE bao phủ phần lớn hành vi lành tính trong khi iForest bắt
các ngoại lệ hiếm.

\subsection{Pipeline SyscallAD}
SyscallAD vận hành trên đặc trưng $174$ chiều tính trên các cửa sổ
syscall cố định (\S\ref{sec:methodology_vi}). VAE không-giám-sát được
huấn luyện trên chuỗi lành tính và chấm điểm cửa sổ kiểm tra bằng lỗi
tái tạo; iForest huấn luyện trên cùng tập lành tính và chấm điểm bằng
độ sâu cô lập. Hai điểm thô được chuẩn hoá và hợp nhất:
$A(W) = \alpha\,A_{\text{VAE}}(W) + (1-\alpha)\,A_{\text{iForest}}(W)$
với $\alpha=0{,}7$.

\subsection{eBPF và Cilium Tetragon}
Tetragon gắn kprobe eBPF vào điểm vào syscall và phơi một API gRPC
dòng (\texttt{FineGuidanceSensors.GetEvents}). Mỗi sự kiện mang định
danh tiến trình và pod, mã syscall, cùng đối số tùy chọn. Tiêu thụ
API trực tiếp được ưu tiên hơn so với phân tích JSON xuất ra vì lược
đồ gRPC có kiểu dữ liệu rõ ràng và thứ tự sự kiện được bảo đảm.

\FloatBarrier
\section{Phương Pháp Luận}
\label{sec:methodology_vi}

\subsection{Trích đặc trưng}
\label{subsec:features_vi}
Mỗi quyết định chấm điểm tính trên cửa sổ $W=200$ định danh syscall
với bước trượt $S=50$. Cửa sổ được ánh xạ về véc-tơ $174$ chiều theo
năm nhóm của thiết kế DeSFAM:
\begin{itemize}
  \item $60$ tần suất syscall (số đếm tương đối của các định danh
        syscall phổ biến nhất trong DongTing);
  \item $40$ cờ nhị phân phân biệt (sự hiện diện của các định danh
        thiên về tấn công);
  \item $8$ thống kê phân phối gồm độ dài cửa sổ, độ đa dạng định
        danh duy nhất, và entropy Shannon của phân phối định danh
        trong cửa sổ;
  \item $40$ tần suất chuyển tiếp bigram (tần suất cặp liền kề phân
        biệt nhất);
  \item one-hot phiên bản nhân $26$ chiều.
\end{itemize}
\texttt{StandardScaler} khớp trên cửa sổ huấn luyện DongTing được áp
dụng nhất quán trong mọi thí nghiệm để bảo đảm so sánh tương đương.

\subsection{Các biến thể mô hình}
\label{subsec:models_vi}
Sáu biến thể mô hình được đánh giá. Siêu tham số được liệt kê ở
Bảng~\ref{tab:v2_hyperparams_vi}. Hậu tố dữ liệu huấn luyện là một
phần của tên biến thể để phép so sánh không nhập nhằng.

\begin{table*}[!t]
\centering
\caption{Siêu tham số của sáu biến thể. Tất cả vận hành trên cùng cửa
sổ $174$ chiều (\S\ref{subsec:features_vi}) với $W=200$, $S=50$ và
StandardScaler khớp trên tập huấn luyện DongTing.}
\label{tab:v2_hyperparams_vi}
\scriptsize
\setlength{\tabcolsep}{4pt}
\renewcommand{\arraystretch}{1.15}
\begin{tabularx}{\textwidth}{@{}p{3.6cm} X X@{}}
\toprule
\textbf{Biến thể} & \textbf{Kiến trúc / Thuật toán} & \textbf{Cấu hình huấn luyện} \\
\midrule
VAE (DongTing)
  & Enc/Dec $174{\rightarrow}128{\rightarrow}64{\rightarrow}24{\rightarrow}64{\rightarrow}128{\rightarrow}174$, SELU, dropout 0,2
  & Chỉ chuỗi lành tính DongTing (5.487); Adam $10^{-3}$; batch 256; $\le$80 epoch; dừng sớm patience 10; 3 seed; MSE+KL; RobustScaler $p_{1}$--$p_{99}$ trên điểm xác nhận \\
VAE (Container)
  & Cùng kiến trúc encoder/decoder
  & Chỉ tập con lành tính Container (149 cửa sổ); cùng vòng huấn luyện; $\le$50 epoch; dừng sớm patience 8; 3 seed \\
iForest (DongTing)
  & Scikit-learn \texttt{IsolationForest}
  & 300 cây; \texttt{contamination='auto'}; \texttt{random\_state=42}; chỉ chuỗi lành tính DongTing; RobustScaler $p_{1}$--$p_{99}$ \\
VAE+iForest (DongTing)
  & $A = \alpha A_{\text{VAE}} + (1{-}\alpha) A_{\text{iForest}}$, cả hai huấn luyện trên DongTing
  & $\alpha = 0{,}7$~\cite{desfam2025}; cả hai điểm chuẩn hoá bằng RobustScaler trước khi hợp nhất \\
LSTM (DongTing+Container)
  & Reshape(174,1) $\to$ LSTM(64) $\to$ Dense(32, ReLU) $\to$ Drop 0,3 $\to$ Dense(1, sigmoid)
  & DongTing+Container có nhãn (cả 2 lớp); BCE; Adam $10^{-3}$; batch 256; $\le$50 epoch; dừng sớm patience 6 theo val-AUC \\
1-D CNN (DongTing+Container)
  & Reshape(174,1) $\to$ Conv1D(32, k=5) $\to$ MaxPool1D(2) $\to$ Conv1D(64, k=3) $\to$ GlobalMaxPool $\to$ Dense(32, ReLU) $\to$ Drop 0,3 $\to$ Dense(1, sigmoid)
  & Cùng cấu hình có-giám-sát với LSTM (DongTing+Container) \\
\bottomrule
\end{tabularx}
\end{table*}

Các biến thể không-giám-sát (VAE, iForest, VAE+iForest) chỉ cần chuỗi
lành tính nên cho phép lựa chọn sạch giữa hai phân phối huấn luyện,
DongTing hoặc Container. Các biến thể có-giám-sát (LSTM, 1-D~CNN) cần
mẫu có nhãn của cả hai lớp; riêng tập tấn công Container chỉ chứa
$\approx\!40$ cửa sổ tấn công, không đủ. Do đó các mô hình có-giám-sát
được huấn luyện trên hợp DongTing+Container.

\subsection{Bộ dữ liệu}
\label{subsec:datasets_vi}

\paragraph{DongTing.}
Bộ dữ liệu đã công bố~\cite{dongting} có $18.966$ chuỗi ($6.850$ lành
tính, $12.116$ độc hại) lấy từ $26$ phiên bản nhân Linux. Chúng tôi
dùng phân chia DTDS công bố ($14.843$ huấn luyện, $1.812$ xác nhận,
$2.311$ kiểm tra) và $5.487$ chuỗi lành tính của tập huấn luyện cho
các phép khớp không-giám-sát.

\paragraph{Container.}
Tập dữ liệu $\approx\!16.200$ syscall được thu trực tiếp từ luồng gRPC
của Cilium Tetragon trên cụm Kubernetes thí nghiệm
(Mục~\ref{sec:deployment_vi}). Thành phần lành tính được bắt từ ba
khối công việc dưới tải thực tế: \texttt{nginx}+\texttt{wrk},
\texttt{redis}+\texttt{redis-benchmark} và
\texttt{postgres}+\texttt{pgbench}. Thành phần tấn công được bắt từ
năm khối công việc: vết khai thác CVE-2021-4034 (Polkit pkexec) $60$
vòng, vết khai thác CVE-2022-0847 (Dirty Pipe) $60$ vòng, một driver
reverse shell, một driver \texttt{tar} exfil, và một mô phỏng miner.
Sau khi trượt cửa sổ ($W=200$, $S=50$) và áp dụng scaler khớp trên
DongTing, tập Container gồm $296$ cửa sổ với tỉ lệ tấn công $28{,}0\%$,
phân chia $70/15/15$ stratified theo khối công việc
(Bảng~\ref{tab:dataset_summary_vi}).

\begin{table}[!t]
\centering
\caption{Phân chia tập dữ liệu Container.}
\label{tab:dataset_summary_vi}
\begin{tabular}{lrr}
\toprule
Phân chia & Cửa sổ & Tỉ lệ tấn công \\
\midrule
huấn luyện  & 207  & 28,0\% \\
xác nhận    & 44   & 27,3\% \\
kiểm tra    & 45   & 28,9\% \\
\bottomrule
\end{tabular}
\end{table}

\subsection{Quy trình đánh giá}
Tất cả biến thể được đánh giá trên cùng hai tập kiểm tra --- kiểm tra
DongTing ($2.311$ cửa sổ, $70{,}7\%$ tấn công) và kiểm tra Container
($45$ cửa sổ, $28{,}9\%$ tấn công). Với mỗi biến thể, ngưỡng quyết
định được chọn theo cực đại Youden-$J$ trên tập xác nhận của phân
phối huấn luyện gốc. Các chỉ số báo cáo gồm AUC, Average Precision
(AP), $F_{1}$, độ chính xác và độ thu hồi.

Với câu hỏi nghiên cứu CH4, bốn vết độc hại (hai khai thác CVE, exfil,
reverse shell) được chấm điểm theo từng cửa sổ. \emph{Độ trễ phát hiện}
của một mô hình trên một vết là số cửa sổ từ đầu vết đến cửa sổ đầu
tiên có điểm vượt ngưỡng của mô hình; mô hình không vượt được báo cáo
là không phát hiện.

\FloatBarrier
\section{Tái Hiện SyscallAD trên DongTing}
\label{sec:reproduction_vi}

\subsection{Lỗi chuẩn hoá trong pipeline tham chiếu}
\label{subsec:normfix_vi}
Bản hiện thực tham chiếu phát hành cùng DongTing chuẩn hoá hai kênh
điểm thô bằng thống kê min--max khớp trên điểm tập huấn luyện. Trên
phân phối DongTing, lỗi tái tạo VAE thô trải qua xấp xỉ sáu cấp độ
độ lớn (giá trị hợp lệ $\approx\!4$, ngoại lệ
$\approx\!1{,}5\times10^{7}$), nên min--max nén phần lớn phân phối
vào một dải mỏng quanh giá trị 0. Ngay cả với $\alpha=0{,}7$, kết
quả tập hợp \emph{suy biến} thành kết quả iForest: cả hai đều cho
$F_{1}=0{,}9151$ với ma trận nhầm lẫn trùng khớp trên tập kiểm tra
DongTing.

Chúng tôi thay min--max bằng dải lượng phân vị khớp trên cùng điểm
huấn luyện:
\begin{equation}
\hat{x} = \mathrm{clip}\!\left(\frac{x - P_{1}(x_{\text{ref}})}
                                    {P_{99}(x_{\text{ref}}) -
                                     P_{1}(x_{\text{ref}}) + \epsilon},
                               0, 1\right).
\label{eq:robust_vi}
\end{equation}
Trên cùng dữ liệu, biên trên của kênh VAE giảm từ
$1{,}5\times10^{7}$ xuống $2{,}18\times10^{3}$, đóng góp của VAE được
bảo toàn trong hợp nhất, và $F_{1}$ tập hợp tăng lên $0{,}9318$. Bài
học có giá trị tổng quát: khi một tập hợp cộng kết hợp điểm từ các
bộ phát hiện không đồng nhất, dải động của từng kênh phải được khảo
sát trước khi chọn bộ chuẩn hoá.

\subsection{Kết quả tái hiện}
Bảng~\ref{tab:repro_vi} báo cáo các chỉ số trên tập kiểm tra DongTing.
$F_{1}$ tập hợp ($0{,}932$) nằm trong khoảng nhiễu lấy mẫu so với
$0{,}92$ đã công bố. Thành phần VAE đơn lẻ đạt $F_{1}=0{,}963$, vượt
tập hợp ở mọi chỉ số tổng hợp trừ độ thu hồi. Đây là chứng cứ đầu
tiên cho thấy công thức tập hợp công bố không thực sự thống trị
thành phần đơn lẻ mạnh nhất của nó.

\begin{table}[!t]
\centering
\caption{Tái hiện trên tập kiểm tra DongTing ($2.311$ cửa sổ,
$70{,}7\%$ tấn công).}
\label{tab:repro_vi}
\begin{tabular}{lrrrrr}
\toprule
Mô hình & AUC & AP & $F_{1}$ & Độ chính xác & Độ thu hồi \\
\midrule
iForest (DongTing)                       & 0,882 & 0,949 & 0,915 & 0,848 & 0,994 \\
VAE (DongTing)                           & 0,960 & 0,970 & 0,963 & 0,937 & 0,990 \\
VAE+iForest (DongTing, $\alpha=0{,}7$)   & 0,921 & 0,960 & 0,932 & 0,876 & 0,996 \\
\midrule
Bài báo~\cite{desfam2025}                & 0,94  & 0,87  & 0,92  & 0,94  & 0,90  \\
\bottomrule
\end{tabular}
\end{table}

Kết quả này trả lời \textbf{CH1} ở dạng khẳng định, có điều kiện là
thay thế bộ chuẩn hoá min--max bị hỏng bằng Phương trình~\ref{eq:robust_vi}.

\FloatBarrier
\section{Triển Khai Qua Cilium Tetragon}
\label{sec:deployment_vi}

\subsection{Kiến trúc}
Bộ phát hiện đã huấn luyện được đóng gói thành một image Python duy
nhất, không phơi cổng và tiêu thụ sự kiện qua dịch vụ gRPC của
Tetragon trong cụm. Với mỗi sự kiện, bộ lọc không gian tên được áp
dụng, định danh syscall được trích từ oneof \texttt{ProcessKprobe},
\texttt{ProcessTracepoint}, \texttt{ProcessExec} hoặc
\texttt{ProcessExit}, và thêm vào hàng đợi cuộn theo từng pod. Mỗi
$S=50$ syscall mới, véc-tơ đặc trưng $174$ chiều đầy đủ được tính lại
và chấm điểm. Phán định được phát ra dưới dạng dòng stdout có cấu
trúc, được hệ thống log của cụm ingest.

\subsection{Ghi nhận triển khai}
Ba mục cấu hình yêu cầu can thiệp đặc thù nền tảng trên Docker Desktop
Kubernetes (nhân $6{,}12$, arm64):

\begin{itemize}
  \item \texttt{/sys/fs/bpf} của máy chủ được mount riêng trên máy ảo
        LinuxKit, làm cho DaemonSet của Tetragon không gắn được dù
        dùng mount propagation Bidirectional hay HostToContainer. Tình
        trạng này được giải quyết bởi một pod đặc-quyền một-lần dùng
        \texttt{nsenter} vào không gian mount của máy chủ và remount
        BPF FS dạng shared.
  \item Mục tiêu kprobe trong \texttt{TracingPolicy} phải dùng dạng
        không tiền tố (\texttt{sys\_execve}) với \texttt{syscall: true}.
        Tetragon cung cấp tiền tố đặc thù kiến trúc khi nạp. Mã hoá
        cứng \texttt{\_\_x64\_sys\_*} khiến chính sách nạp im lặng dưới
        dạng no-op trên arm64.
  \item Helm chart binding máy chủ gRPC vào unix socket mặc định. Đặt
        \texttt{tetragon.grpc.address=0.0.0.0:54321} và phơi Service
        \texttt{ClusterIP} trên cổng đó là bắt buộc để API tiếp cận
        được từ pod khác.
\end{itemize}

\subsection{Đặc tính hiệu năng}
Độ trễ chấm điểm theo cửa sổ, đo end-to-end từ khi sự kiện đến qua
xây dựng đặc trưng và suy luận mô hình, nằm trong $9$--$18$~ms cho
mọi biến thể trên một lõi CPU đơn. Năm lần lấy mẫu được dùng cho
mỗi dự đoán VAE để ổn định ước lượng lỗi tái tạo. Các biến thể
có-giám-sát chậm hơn không đáng kể ($24$--$34$~ms) do bước convolution
thứ hai.

\FloatBarrier
\section{Đánh Giá Thực Nghiệm}
\label{sec:evaluation_vi}

\subsection{Kiểm thử đa dạng sáu khối công việc (v1)}
\label{subsec:diversity_vi}
Bộ phát hiện v1 đã triển khai với tập hợp VAE+iForest~(DongTing) gốc
được chạy với sáu khối công việc trong không gian tên \texttt{demo}:
một bomb fork-exec, một bộ quét mạng, vòng lặp I/O \texttt{dd}/\texttt{cat},
vòng lặp leo thang đặc quyền \texttt{unshare}+\texttt{mount}, vòng lặp
gắn trình gỡ lỗi, và đối chứng lành tính \texttt{curl}.
Bảng~\ref{tab:diversity_vi} tổng hợp điểm theo cửa sổ quan sát được.

\begin{table}[!t]
\centering
\caption{Phán định v1 trên sáu khối công việc đa dạng với
VAE+iForest~(DongTing), $\alpha=0{,}7$, ngưỡng $0{,}077$.}
\label{tab:diversity_vi}
\begin{tabular}{lrrrr}
\toprule
Khối công việc          & $n$ & TB VAE & TB iForest & TB kết hợp \\
\midrule
\texttt{anom-execbomb}  & 69  & \textbf{0,524} & 0,523 & \textbf{0,524} \\
\texttt{anom-netscan}   & 7   & 0,408          & 0,594 & 0,464          \\
\texttt{anom-kernel}    & 24  & 0,274          & 0,485 & 0,337          \\
\texttt{anom-privesc}   & 228 & 0,222          & 0,536 & 0,316          \\
\texttt{anom-benign}    & 47  & 0,170          & 0,514 & 0,273          \\
\texttt{anom-ptrace}    & 241 & \textbf{0,076} & 0,420 & \textbf{0,179} \\
\bottomrule
\end{tabular}
\end{table}

Hai mẫu hiện rõ trong Bảng~\ref{tab:diversity_vi} trả lời \textbf{CH2}
ở dạng phủ định. Thứ nhất, đối chứng lành tính được chấm điểm trên
ngưỡng công bố bởi mọi thành phần, cho thấy ngưỡng tinh chỉnh theo
DongTing không chuyển sang đường nền container thật. Thứ hai, khối
gắn trình gỡ lỗi (sử dụng \texttt{ptrace} dày đặc) được chấm điểm
\emph{thấp hơn} đối chứng lành tính, không phải cao hơn. Chúng tôi
diễn giải đây là hệ quả của việc VAE không-giám-sát đã học, từ lớp
lành tính bộ kiểm thử nhân của DongTing, rằng lưu lượng \texttt{ptrace}
dày đặc là bình thường --- vì lớp đó thực tế bão hoà luồng làm việc
gỡ lỗi.

\subsection{Phân tích so sánh mô hình}
\label{subsec:comparative_vi}
Bảng~\ref{tab:v2_models_vi} báo cáo các chỉ số của sáu biến thể trên
cả tập kiểm tra DongTing và tập kiểm tra Container.

\begin{table}[!t]
\centering
\caption{Hiệu năng theo từng tập kiểm tra. DT~=~kiểm tra DongTing
($2.311$ cửa sổ, $70{,}7\%$ tấn công); Cont.~=~kiểm tra Container
($45$ cửa sổ, $28{,}9\%$ tấn công).}
\label{tab:v2_models_vi}
\begin{tabular}{lrrrr}
\toprule
Mô hình                           & AUC DT & AUC Cont. & $F_{1}$ DT & $F_{1}$ Cont. \\
\midrule
VAE (DongTing)                    & 0,960  & \textbf{0,135} & 0,952 & 0,481 \\
VAE (Container)                   & 0,687  & 1,000          & 0,842 & 1,000 \\
iForest (DongTing)                & 0,882  & 1,000          & 0,832 & 1,000 \\
LSTM (DongTing+Container)         & 0,998  & 1,000          & 0,992 & 1,000 \\
1-D CNN (DongTing+Container)      & 1,000  & 1,000          & 1,000 & 1,000 \\
\bottomrule
\end{tabular}
\end{table}

VAE~(DongTing) sụp đổ trên tập kiểm tra Container với AUC$\,=\,0{,}135$,
thấp hơn ngẫu nhiên về mặt thống kê. Huấn luyện lại cùng kiến trúc
trên tập con lành tính Container đưa AUC Container về $1{,}000$. Các
baseline có-giám-sát đạt chỉ số tổng hợp mạnh nhất, nhưng kèm lưu ý
nêu ở Mục~\ref{subsec:models_vi}: khoảng $70\%$ số cửa sổ tấn công
Container đã được thấy trong huấn luyện, làm tăng các chỉ số Container
so với một đánh giá ngoài-mẫu thực sự.

Bằng chứng kết hợp từ Bảng~\ref{tab:repro_vi} và
Bảng~\ref{tab:v2_models_vi} trả lời \textbf{CH3}: trên không gian
đặc trưng có sẵn dữ liệu tấn công có nhãn, các bộ phân loại có-giám-sát
đơn giản (LSTM, 1-D~CNN) ngang hoặc vượt tập hợp không-giám-sát đã
công bố. Thành phần iForest đóng góp tối thiểu; điểm của cột iForest
trong Bảng~\ref{tab:diversity_vi} gần như không đổi theo khối công
việc, cho thấy tập hợp thừa kế phần lớn sức mạnh phân biệt từ VAE.

\subsection{Độ trễ phát hiện CVE}
\label{subsec:cve_latency_vi}
Bảng~\ref{tab:cve_latency_vi} báo cáo độ trễ phát hiện (theo số cửa
sổ) cho bốn vết độc hại của Mục~\ref{subsec:datasets_vi}.

\begin{table}[!t]
\centering
\caption{Độ trễ phát hiện theo cửa sổ (`--' $=$ không phát hiện trong
vết). DT $=$ huấn luyện DongTing, Cont. $=$ huấn luyện Container,
$^{\dag}$ $=$ mô hình có-giám-sát huấn luyện trên DongTing+Container.}
\label{tab:cve_latency_vi}
\begin{tabular}{lrrrrrr}
\toprule
Vết            & $n_w$ & VAE\,(DT) & VAE\,(Cont.) & iForest\,(DT) & LSTM\,$^{\dag}$ & CNN\,$^{\dag}$ \\
\midrule
CVE-2021-4034  & 63    & 0         & --           & 0             & 4               & 0              \\
CVE-2022-0847  & 8     & 0         & 0            & 0             & --              & 0              \\
exfil          & 7     & 0         & 0            & 0             & --              & 0              \\
revshell       & 5     & 0         & 4            & 0             & --              & 0              \\
\bottomrule
\end{tabular}
\end{table}

Các biến thể huấn luyện DongTing và iForest đánh dấu mọi vết độc hại
ngay cửa sổ~$0$, nhưng để diễn giải cần đối chiếu với đường nền đa
dạng ở Mục~\ref{subsec:diversity_vi}: các biến thể này cũng đánh dấu
đối chứng lành tính. Độ trễ bằng 0 chỉ là phụ phẩm của việc gán nhãn
dương không phân biệt chứ không phải phát hiện phân biệt. LSTM
có-giám-sát bỏ sót ba trên bốn vết. 1-D~CNN có-giám-sát phát hiện cả
bốn ngay cửa sổ~$0$ với điểm cực đại chuẩn hoá nằm trong $[0{,}95,
0{,}99]$. Đây là câu trả lời vận hành cho \textbf{CH4}: khi mô hình
có-giám-sát đã thấy các ví dụ tấn công tương đương trong huấn luyện,
cửa sổ chấm điểm đầu tiên là đủ.

\subsection{Triển khai lại sản xuất}
\label{subsec:redeploy_vi}
Để xác nhận sự cải thiện về chất trong triển khai, image bộ phát hiện
được rebuild với cờ \texttt{--model-variant} và triển khai lại với
VAE~(Container) cùng ngưỡng quyết định $0{,}5$ chọn từ phân phối điểm
v2. Bài kiểm thử cùng sáu khối công việc được lặp lại.

\begin{table}[!t]
\centering
\caption{Kiểm thử triển khai sáu khối công việc, trước và sau.
Trước: VAE+iForest~(DongTing) ngưỡng $0{,}077$.
Sau: VAE~(Container) ngưỡng $0{,}5$.}
\label{tab:v2_multipod_vi}
\begin{tabular}{lrrr}
\toprule
Khối công việc         & TB trước & TB sau & ATK\,\% sau \\
\midrule
\texttt{anom-netscan}   & 0,464 & 1,000 & 100\% \\
\texttt{anom-kernel}    & 0,337 & 1,000 & 100\% \\
\texttt{anom-execbomb}  & 0,524 & 1,000 & 100\% \\
\texttt{anom-privesc}   & 0,316 & 0,867 & 85\%  \\
\texttt{anom-ptrace}    & 0,179 & 0,303 & 0\%   \\
\texttt{anom-benign}    & 0,273 & 0,218 & 0\%   \\
\bottomrule
\end{tabular}
\end{table}

Bốn dạng lỗi đã nhận diện trước đó được giải quyết. Đối chứng lành
tính không còn trên ngưỡng; điểm của khối \texttt{ptrace} nay vượt
điểm đối chứng lành tính; biên độ giữa tấn công xấu nhất và lành
tính tốt nhất nới từ $\approx\!3\times$ ở v1 lên $\approx\!4{,}6\times$
ở v2; và bốn bất thường thật được gắn nhãn ATTACK với tỉ lệ $\ge\!85\%$
trong khi cả hai khối lớp lành tính vẫn yên tĩnh.

\FloatBarrier
\section{Thảo Luận}
\label{sec:discussion_vi}

Kết quả chính là việc lựa chọn dữ liệu huấn luyện chi phối lựa chọn
kiến trúc trong lớp bài toán này. Giữ nguyên không gian đặc trưng
$174$ chiều, kiến trúc VAE và bộ chuẩn hoá RobustScaler, huấn luyện
lại chỉ trên dữ liệu Container đưa AUC Container từ dưới ngẫu nhiên
lên $1{,}0$. Lớp lành tính của DongTing không đại diện cho mẫu syscall
container sản xuất vì nó được tạo bởi các bộ kiểm thử nhân vận hành
các syscall hiếm dùng --- bao gồm các luồng \texttt{ptrace} dày đặc
làm ô nhiễm phân phối lành tính đã học. Người triển khai SyscallAD
trên khối công việc thật do đó nên dự liệu một pha huấn luyện
trong-môi-trường thay vì dựa vào trọng số huấn luyện trên DongTing.

Một kết quả phụ liên quan đến cấu thành tập hợp. Cả ở phép tái hiện
trên kiểm tra DongTing (Bảng~\ref{tab:repro_vi}) và phép kiểm thử đa
dạng ngoài phân phối (Bảng~\ref{tab:diversity_vi}), thành phần VAE
đơn lẻ vượt tập hợp $\alpha=0{,}7$ VAE+iForest ở các chỉ số tổng hợp.
Điểm của iForest theo cửa sổ nằm trong dải hẹp $[0{,}42, 0{,}59]$,
đóng góp ít tín hiệu phân biệt. Chúng tôi giả thiết rằng lựa chọn
$\alpha=0{,}7$ của bài báo chỉ đúng dưới một bộ chuẩn hoá không được
nêu rõ, và dưới bộ chuẩn hoá robust nguyên tắc của
Phương trình~\ref{eq:robust_vi}, chế độ $\alpha\!\to\!1$ được ưu tiên.

\paragraph{Các nguy cơ với tính hợp lệ.}
Thứ nhất, tập dữ liệu Container nhỏ ($296$ cửa sổ, $\approx\!16$k
syscall). Các baseline có-giám-sát hưởng lợi không cân xứng vì quan
sát hầu hết phân phối tấn công trong huấn luyện; các kết quả kiểm tra
Container của chúng do đó nên đọc như giới hạn trên. Một kho tấn công
Container lớn hơn sẽ thắt chặt phép so sánh. Thứ hai, TracingPolicy
chỉ hook một tập tuyển chọn $20$ syscall, làm nén tỉ lệ sự kiện thực
tế theo từng pod xuống $\approx\!10$ sự kiện/giây và thiên lệch phân
phối được bắt về phía các nguyên hàm điều khiển luồng và tạo tiến
trình. Một chính sách raw-tracepoint sẽ cho vết dày hơn với khối
lượng sự kiện cao hơn. Thứ ba, mọi phép đo được thực hiện trên một
node Docker Desktop đơn (LinuxKit, arm64); các con số độ trễ tuyệt
đối không nhất thiết chuyển sang cụm nhiều node.

\FloatBarrier
\section{Giới Hạn và Công Việc Tương Lai}
\label{sec:future_vi}

Tập dữ liệu Container đủ để bác bỏ mô hình huấn luyện trên DongTing
trong triển khai nhưng không đủ để huấn luyện một baseline có-giám-sát
mạnh chỉ trên Container. Bước tự nhiên tiếp theo là mở rộng pipeline
thu thập dữ liệu sang nhiều node và nhiều ngày lưu lượng. Với cơ sở
đó, bốn thí nghiệm tiếp theo gợi ý sẵn sàng: (i)~thay VAE+iForest
bằng chỉ-VAE ở ngưỡng triển khai-thời để kiểm chứng liệu giá trị
biên của tập hợp có biện minh được chi phí; (ii)~quét
$\alpha\in\{0{,}7; 0{,}8; 0{,}9; 1{,}0\}$ dưới
Phương trình~\ref{eq:robust_vi}; (iii)~huấn luyện một bộ phân loại
có-giám-sát chỉ trên dữ liệu Container khi kho tấn công đủ lớn, rồi
so sánh với VAE~(Container); (iv)~ghép cảnh báo với hành động thực
thi \texttt{Sigkill} của Tetragon để đo tỉ lệ ngăn chặn end-to-end
trước khai thác CVE, không chỉ phát hiện.

\FloatBarrier
\section{Kết Luận}
\label{sec:conclusion_vi}

Chúng tôi đã tái hiện SyscallAD trên bộ dữ liệu công bố trong phạm
vi nhiễu lấy mẫu của các chỉ số tiêu đề, có điều kiện là thay một
bộ chuẩn hoá không được nêu bằng một bộ scaler dải lượng phân vị
nguyên tắc. Sau đó chúng tôi triển khai bộ phát hiện trên luồng sự
kiện gRPC của Cilium Tetragon và quan sát thấy mô hình đã công bố
không chuyển giao được sang các khối công việc container kiểu sản
xuất: một vòng lặp \texttt{curl} lành tính bị gắn nhãn còn vòng lặp
gắn trình gỡ lỗi thì không. Huấn luyện lại cùng kiến trúc VAE trên
vết syscall thu từ khối công việc container thật loại bỏ sạch lỗi
này. Đóng góp của thành phần iForest đối với khả năng phân biệt
được chứng minh là nhỏ trong cả hai cấu hình, và các bộ phân loại
có-giám-sát trên cùng không gian đặc trưng cạnh tranh được. Khuyến
nghị thực tiễn là mọi triển khai bộ phát hiện kiểu SyscallAD nên đi
trước bởi việc thu thập dữ liệu trong-môi-trường, và bộ chuẩn hoá
cùng ngưỡng quyết định nên được khảo sát trên phân phối điểm của
môi-trường-triển-khai thay vì kế thừa từ mặc định của
môi-trường-huấn-luyện.
